\specialchapter{历史注记}

（第 I 至 V 章）

特征值与傅里叶级数的概念在 1830 年已为人熟知，并贯穿十九世纪的许多分析领域（参见~ÉHM, pp.~114, 260, 275）。但谱理论与调和分析——按我们在这本书中所理解的意义——直到大约一个世纪之后才开始具有其现在的形式；这与它们在 1930 年前后同群表示研究的汇合、随后在 1940 年前后同赋范代数研究的汇合是同时发生的。

我们在这里仅限于勾画它们从 1906 年（希尔伯特在其积分方程工作中引入“连续谱”概念之时）到 1945–1950 年间演化的主要脉络；在此期间，这里所陈述的理论基本上获得了保持至今的形式。

